\chapter{Аналитический раздел}

В данном разделе анализируется предметная область, существующие решения,
формализуются требования к разрабатываемому комплексу.

\section{Анализ предметной области}

Наиболее часто встречающимся интерфейсом для подключения внешних запоминающих
устройств к вычислительной технике является USB (Universal Serial Bus). К таким устройствам
относятся флеш-карты и жёсткие диски. Наиболее простой способ устранения рисков, связанных с 
использованием внешних запоминающих устройств, --- программное отключение USB-портов.
Однако это не всегда может быть удобно, так как порты используются и устройствами ввода-вывода, 
например, мышью или клавиатурой. Более рациональным решением будет не полностью отключать порты, 
а программно изменять поведение компьютера при подключении устройства через USB-порт.

\section{Сравнительный анализ существующих решений}

Все существующие продукты делятся на две группы --- корпоративные системы предотвращения
утечек и утери данных, и решения для личного пользования. Корпоративные решения имеют 
распределённую архитектуру, рассчитаны на управление и мониторинг множества рабочих станций, и являются
избыточными для персонального компьютера. В анализе участвуют только программы для личного
использования.

Для анализа были выбраны следующие программные комплексы:

\begin{itemize}[label=---]
    \item USBGuard~\cite{USBGuard};
    \item USBFilter~\cite{USBFilter};
    \item USBWall~\cite{USBWall};
    \item USBShield~\cite{USBShield};
    \item Ratool~\cite{Ratool};
    \item GiliSoft USB Lock~\cite{GilisoftUSBLock};
    \item USB Block~\cite{USBBlock}.
\end{itemize}

Анализ проводится по следующим критериям:

\begin{enumerate}
    \item поддержка ОС (операционной системы) Linux Ubuntu;
    \item возможность создания списка доверенных устройств;
    \item возможность установки даты истечения периода доверия к устройству;
    \item наличие графического интерфейса;
    \item наличие обновлений и поддержки.
\end{enumerate}

В таблице~\ref{tab:compare_solutions} представлены результаты анализа существующих решений.

\begin{table}[H]
    \caption{Результат анализа существующих решений}
    \begin{tabular}{|l|c|c|c|c|c|c|c|}
    \hline
    \diagbox{\textbf{Номер критерия}}{\textbf{Номер решения}} & \textbf{1} & \textbf{2} & \textbf{3} & \textbf{4} & \textbf{5} & \textbf{6} & \textbf{7}\\
    \hline
    1 & + & + & + & - & - & - & - \\ \hline
    2 & + & + & + & + & - & + & + \\ \hline
    3 & - & - & - & - & - & - & - \\ \hline
    4 & - & - & - & + & + & + & + \\ \hline
    5 & + & - & - & + & + & + & + \\ \hline
    \end{tabular}
    \label{tab:compare_solutions}
\end{table}

Таким образом, ни одно из представленных решений не соответствует в полной мере поставленным критериям.

\section{Диаграмма прецедентов разрабатываемого комплекса}

На рисунке~\ref{img:use_cases} представлена диаграмма прецедентов разрабатываемого комплекса.

\includeimage
    {use_cases}
    {f}
    {H}
    {0.70\textwidth}
    {Диаграмма прецедентов комплекса}

\section{Формализация и описание информации, подлежащей хранению в списке доверенных устройств}

Формальное описание доверенного устройства может быть представлено следующими аттрибутами:

\begin{enumerate}
    \item идентификатор производителя;
    \item идентификатор устройства;
    \item уникальный серийный номер устройства;
    \item наименование производителя;
    \item наименование устройства;
    \item дата истечения доверия.
\end{enumerate}

Уникальность устройства должна определять комбинация серийного номера и идентификаторов
производителя и устройства.

\section{Анализ решений для хранения списка доверенных устройств}

Существует несколько основных способов организации хранилища списка доверенных устройств, 
к которым относятся базы данных и плоские файлы. Анализ проводится по следующим критериям:

\begin{enumerate}
    \item наличие встроенных механизмов поддержания целостности данных (ограничения, уникальность, проверки);
    \item устойчивость хранилища к сбоям и некорректному завершению работы (восстановление согласованного состояния);
    \item отсутствие дополнительных зависимостей у прикладного приложения;
    \item отсутствие возможности несанкционированного ручного просмотра и редактирования данных без специализированных инструментов;
    \item возможность выполнения выборок и фильтрации данных без дополнительной реализации алгоритмов поиска;
    \item отсутствие необходимости инициализации и настройки хранилища перед началом работы;
    \item возможность расширения структуры данных и добавления новых требований без переработки существующей логики хранения.
\end{enumerate}

В таблице~\ref{tab:compare_storages} представлены результаты анализа хранилищ данных.


\begin{table}[H]
    \caption{Результат анализа решений для хранения списка доверенных устройств}
    \begin{tabular}{|l|c|c|c|c|c|c|c|}
    \hline
    \diagbox{\textbf{Номер критерия}}{\textbf{Решение}} & \textbf{Плоские файлы} & \textbf{Базы данных}\\
    \hline
    1 & - & + \\ \hline
    2 & - & + \\ \hline
    3 & + & - \\ \hline
    4 & - & + \\ \hline
    5 & - & + \\ \hline
    6 & + & - \\ \hline
    7 & - & + \\ \hline
    \end{tabular}
    \label{tab:compare_storages}
\end{table}

На основе таблицы анализа для разрабатываемого комплекса предпочтительнее использовать базу данных.

\section{Анализ существующих баз данных на основе формализации данных}

Существует три основных типа баз данных:

\begin{enumerate}
    \item дореляционные;
    \item реляционные;
    \item постреляционные.
\end{enumerate}

\subsection{Дореляционные базы данных}
Дореляционные базы данных~\cite{DBTypes} хранят записи в виде связей предок/потомок. 
Основными недостатками дореляционных баз данных являются избыточность данных и 
сложность выполнения запросов. Изменение схемы базы данных крайне ограничена. 
Целостность данных обеспечивается за счёт жёсткой структуры связей.

\subsection{Реляционные базы данных}
В реляционных базах данных данные хранятся в виде двумерных таблиц.
Связь между ними обеспечивается общими столбцами, такими как идентификаторы. 
Табличная архитектура подходит для хранения структурированной информации.
Целостность и непротиворечивость данных соблюдается при помощи ограничений.

\subsection{Постреляционные базы данных}
Постреляционные базы данных подходят для хранения неструктурированных
данных, поддерживая логические связи между ними. Недостатком является сложность 
обеспечения целостности данных. Отличаются хорошей горизонтальной
масштабируемостью, сложность запросов варьируется в зависимости от конкретной реализации.

\subsection{Вывод}
Для реализации поставленной задачи была выбрана реляционная модель, так как она лучше всего 
подходит для хранения структурированной информации, обеспечивая баланс между целостностью данных, 
простотой запросов и возможностями масштабирования.

\section{Выводы}

В данном разделе была проанализирована предметная область, проведено сравнение существующих решений
по сформулированным критериям. Программно-алгоритмический комплекс был формализован с использованием
диаграммы прецедентов, описано представление доверенного устройства, проанализированы способы
организации хранилища списка, выбран наиболее подходящий тип базы данных.